\documentclass[11pt]{article}
\usepackage[a4paper,margin=2cm]{geometry}
\usepackage[T1]{fontenc}
\usepackage[utf8]{inputenc}
\usepackage{amsmath,amssymb,longtable,array,booktabs}
\usepackage{ragged2e}
\renewcommand{\arraystretch}{1.25}
\newcolumntype{L}[1]{>{\raggedright\arraybackslash}p{#1}}

\title{Glosario de funciones\\
\large ``Bound states and resonances of screened Coulomb potentials with
the FEAST contour-integration eigensolver in an orthonormal Laguerre
basis''}
\author{Documento de referencia r\'apida --- reuni\'on con R. Manzanilla}
\date{\today}

\begin{document}
\maketitle
\noindent Toda funci\'on matem\'atica y toda rutina num\'erica que aparece
en el manuscrito o que la sustenta, con su nombre expl\'icito. Este
documento responde directamente al comentario telef\'onico del profesor
Manzanilla: ``usamos una funci\'on que no hab\'iamos nombrado como tal''
(identificada como la funci\'on Gamma de Euler) --- ver la nota final.

\begingroup
\footnotesize
\begin{longtable}{@{}L{4.4cm}L{9.2cm}L{2.8cm}@{}}
\toprule
Funci\'on & Qu\'e es y d\'onde se usa & Referencia \\
\midrule
\endhead
Polinomio de Laguerre generalizado $L_j^{(\alpha)}(x)$ & Polinomio
  ortogonal cl\'asico; con $\alpha=2\ell+2$ y el prefactor
  $x^{\ell+1}e^{-x/2}$ construye cada funci\'on base $\varphi_j(r)$ de la
  Ec.~(3). Es \emph{la} funci\'on especial central del m\'etodo. &
  Ec.~(3), Sec.~2.2 \\
Factoriales $j!$, $(j+2\ell+2)!$ & Normalizan la base:
  $\varphi_j=\sqrt{2\beta}\sqrt{j!/(j+2\ell+2)!}\,x^{\ell+1}e^{-x/2}L_j^{(2\ell+2)}(x)$.
  Se escriben expl\'icitamente como factoriales, \textbf{no} como funci\'on
  Gamma $\Gamma(\cdot)$, precisamente para no chocar con el s\'imbolo
  $\Gamma$ de la anchura de resonancia. & Ec.~(3) \\
Recurrencia de tres t\'erminos de los polinomios de Laguerre
  ortonormalizados & Relaci\'on recursiva usada para evaluar $\varphi_j(r)$
  de forma num\'ericamente estable, en vez de la f\'ormula cerrada
  (inestable para $N$ grande). & Sec.~2.2, texto tras Ec.~(3) \\
Identidad de derivada
  $x\,dL_j^{(\alpha)}/dx = jL_j^{(\alpha)}-(j+\alpha)L_{j-1}^{(\alpha)}$ &
  Identidad algebraica usada para construir el t\'ermino cin\'etico de la
  matriz Hamiltoniana sin derivar num\'ericamente. & Sec.~2.3 \\
Factor de apantallamiento $e^{-\lambda r}$ (y su versi\'on rotada
  $e^{-\lambda r e^{i\theta}}$) & Funci\'on exponencial que codifica el
  apantallamiento de Yukawa/Debye en el potencial; se punt\'ua directamente
  dentro de la cuadratura. & Ec.~(2), Sec.~2.3 \\
Transformaci\'on de escalamiento complejo $r\to re^{i\theta}$ & No es una
  ``funci\'on especial'' con nombre propio sino un cambio de variable
  (mapa) que rota el eje radial al plano complejo; deja fijas las
  energ\'ias ligadas y expone las resonancias como autovalores aislados. &
  Ec.~(2), Sec.~2.1 \\
Resolvente $(z-H)^{-1}$ & Funci\'on matricial del par\'ametro complejo $z$;
  es el integrando del proyector espectral de FEAST. & Ec.~(5) \\
Proyector espectral (integral de contorno)
  $P=\frac{1}{2\pi i}\oint_{\mathcal C}(z-H)^{-1}dz$ & Proyecta sobre el
  subespacio de autovalores encerrados por el contorno $\mathcal C$; es el
  objeto matem\'atico que define a FEAST. & Ec.~(5), Sec.~2.4 \\
Cuadratura de Gauss--Legendre & Regla de integraci\'on num\'erica, usada en
  \textbf{dos} lugares distintos del m\'etodo: (a) para evaluar los
  elementos de matriz (integrales de polinomios contra $e^{-x}$, Sec.~2.3);
  (b) para aproximar la integral de contorno de FEAST con $n_c$ nodos $z_k$
  y pesos $\omega_k$ (Sec.~2.4). & Sec.~2.3, Sec.~2.4 \\
Delta de Kronecker $\delta_{ij}$ & Vale 1 si $i=j$, 0 si no; aparece en la
  identidad de ortonormalidad certificada num\'ericamente (Ec.~4). &
  Ec.~(4) \\
\texttt{feast\_interval(\ldots)} & Implementaci\'on propia (Python, c\'odigo
  liberado) del algoritmo FEAST original de Polizzi (2009) sobre un
  \textbf{intervalo real} del eje real; usada para validar contra el
  hidr\'ogeno y contra LAPACK. \textbf{No} se llama a ninguna librer\'ia
  FEAST externa. & Sec.~2.4, \texttt{src/feast.py} \\
\texttt{feast\_region(\ldots)} & Implementaci\'on propia: adaptaci\'on del
  FEAST no-herm\'itico de Kestyn~\emph{et al.}~(2016) con un contorno
  \textbf{circular} centrado en la posici\'on anticipada de la resonancia
  $E_r-i\Gamma/2$. \textbf{Este es el aporte metodol\'ogico del paper} ---
  el punto que el profesor pidi\'o resaltar expl\'icitamente. & Sec.~2.4,
  \texttt{src/feast.py} \\
Paso de Rayleigh--Ritz & Procedimiento est\'andar de \'algebra lineal:
  resuelve el problema de autovalores reducido dentro del subespacio
  proyectado $\mathrm{span}(Q)$ que entrega FEAST. & Sec.~2.4 \\
Diagonalizaci\'on directa (LAPACK) & Diagonalizaci\'on exacta de la misma
  matriz, usada \'unicamente como referencia de validaci\'on cruzada (no es
  parte del m\'etodo propuesto). & Sec.~3 \\
Solver de diferencias finitas (independiente) & Soluci\'on num\'erica del
  problema radial en una malla uniforme; no comparte c\'odigo con la base
  de Laguerre. Se usa solo para validar la Tabla~1. & Sec.~3, Tabla~1 \\
\bottomrule
\end{longtable}
\endgroup

\vspace{1em}
\noindent\textbf{Sobre la funci\'on sin nombrar (aclarado):} en un borrador
anterior la normalizaci\'on de la base se escrib\'ia usando la funci\'on
Gamma de Euler $\Gamma(j+2\ell+3)$, que nunca se nombraba expl\'icitamente
como tal \emph{y} colisionaba notacionalmente con el s\'imbolo $\Gamma$ ya
usado para la anchura de resonancia. La correcci\'on aplicada reescribe esa
normalizaci\'on con el factorial equivalente $(j+2\ell+2)!$
(pues $\Gamma(n{+}1)=n!$ para $n$ entero), eliminando la funci\'on sin
nombrar y la colisi\'on de s\'imbolos a la vez. En el \texttt{main.tex}
actual no aparece ninguna funci\'on Gamma de Euler.

\end{document}
